\section{PHÂN CỤM (CLUSTERING)}

\subsection{Khái niệm \& Mục đích}
Phân chia dữ liệu vào các nhóm (cụm) sao cho các đối tượng trong cùng một cụm tương đồng nhau và khác biệt với các đối tượng ở cụm khác. Mục đích là khám phá cấu trúc tự nhiên của dữ liệu.

\subsection{Khác nhau giữa Phân cụm và Phân lớp}
\begin{table}[H]
    \centering
    \renewcommand{\arraystretch}{1.3}
    \begin{tabularx}{\textwidth}{|p{3cm}|X|X|}
        \hline
        \textbf{Tiêu chí} & \textbf{Phân lớp (Classification)} & \textbf{Phân cụm (Clustering)} \\
        \hline
        \textbf{Loại học} & Có giám sát (Supervised). & Không giám sát (Unsupervised). \\
        \hline
        \textbf{Dữ liệu Input} & Có nhãn (Label). & Không có nhãn. \\
        \hline
        \textbf{Mục đích} & Dự đoán nhãn cho dữ liệu mới. & Tìm cấu trúc/nhóm trong dữ liệu. \\
        \hline
    \end{tabularx}
\end{table}